\documentclass[11pt]{article}
\usepackage{amsmath,amssymb,amsthm,hyperref}
\newtheorem{proposition}{Proposition}
\begin{document}
\title{What dimension counting in first-order parents implies}
\author{}\date{}\maketitle

Let $\mathcal C\subset\operatorname{Sym}^N\mathbb C^2$ be a
$k$-dimensional second-order anticoherent subspace, and let $P$ be its
orthogonal projector. Put $j=N/2$, $c=j(j+1)/3$, and $D=N+1$.
The polarized identities are
\[
 PJ_aP=0,\qquad PJ_aJ_bP=c\delta_{ab}P.
\]
Hence $\mathcal C,J_x\mathcal C,J_y\mathcal C,J_z\mathcal C$ are
pairwise orthogonal $k$-dimensional spaces.

\begin{proposition}[Conditional parent-dimension bound]
If a first-order AC parent $\mathcal W$ contains $\mathcal C$, and
$d=\dim\mathcal W$, then
\[
 \boxed{d+3k\le N+1.}
\]
\end{proposition}
\begin{proof}
For any $w\in\mathcal W$ and $u\in\mathcal C\subseteq\mathcal W$,
first-order anticoherence gives $\langle w,J_a u\rangle=0$.
Thus $\mathcal W$ is orthogonal to each of the three pairwise orthogonal
spaces $J_a\mathcal C$. Their direct sum has dimension $d+3k$ and lies
in the $N+1$ dimensional physical space.
\end{proof}

\paragraph{Why the unrestricted first-order maximum cannot be substituted.}
For $N\ge4$ the established maximum is
$K(N,1)=\lfloor N/3\rfloor+1$. This is an \emph{upper} bound on $d$;
it does not imply that any given $\mathcal C$ has a parent of this
maximum dimension. Substituting $d=K(N,1)$ into the proposition without
an extension theorem reverses the required direction of information.

There is an exact counterexample in the catalogue. Its optimal
four-dimensional second-order code at $N=18$ has a largest first-order
parent of dimension five, although $K(18,1)=7$. The proof and independent
exact verifiers are in
\href{../t2-nested-t1/extension-obstruction.tex}{\texttt{../t2-nested-t1/extension-obstruction.tex}}.
Thus searching every maximum-dimensional first-order parent would still
miss this particular optimal second-order code.

\paragraph{What the hypothetical maximum-parent assumption would give.}
If maximum-dimensional containment is established for a particular code,
then the conditional conclusion is
\[
 k\le\left\lfloor\frac{N-\lfloor N/3\rfloor}{3}\right\rfloor
 =\frac{2N}{9}+O(1).
\]
It can improve a finite-size bound under that extra assumption; for
example at $N=30$ it gives $k\le6$. It is not an unrestricted exclusion
of dimension seven. Asymptotically its coefficient $2/9$ is larger than
$1/5$, so even maximum-parent containment would not improve the leading
coefficient of the already established universal bound
\[
 k\le\min\left\{\left\lfloor\frac{N-2}{4}\right\rfloor,
                   \left\lfloor\frac{N+5}{5}\right\rfloor\right\}
 \qquad(N\ge8).
\]
Using only the always valid $d\ge k$ recovers $4k\le N+1$.
Even a separately justified one-dimensional extension, $d\ge k+1$,
would give only $4k\le N$, weaker than the known $4k\le N-2$.

\paragraph{The exact extension problem retains operator information.}
Define
\[
 R=I-P-\frac1c\sum_aJ_aPJ_a.
\]
It is an orthogonal projector of rank $r=N+1-4k$. The first-order parents
of $\mathcal C$ are precisely
\[
 \mathcal W=\mathcal C\oplus\mathcal A,\qquad
 \mathcal A\subseteq\operatorname{ran}R,\qquad
 P_{\mathcal A}J_aP_{\mathcal A}=0.
\]
Thus determining the largest parent requires solving the common zero
compression problem for the three matrices $RJ_aR$ on the complement.
The rank of $R$ alone does not solve that problem.

There are additional exact trace identities:
\[
 \operatorname{Tr}(RJ_a)=0,\qquad
 \operatorname{Tr}(RJ_aJ_b)=cr\delta_{ab}.
\]
For completeness, set $C=j(j+1)$ and
$\Phi(X)=\sum_aJ_aXJ_a$. Then
\[
 \Phi(J_b)=(C-1)J_b,
\]
\[
 \Phi\!\left(\frac{J_aJ_b+J_bJ_a}{2}\right)
 =(C-3)\frac{J_aJ_b+J_bJ_a}{2}+C\delta_{ab}I.
\]
Self-adjointness of $\Phi$ in the Hilbert--Schmidt pairing, together with
the AC2 identities for $P$, gives the displayed traces; the antisymmetric
part vanishes by the first trace identity. If $r>0$, the density operator
$R/r$ therefore has isotropic first and second spin moments. This is a
mixed-state trace statement. It does not say that every vector in
$\operatorname{ran}R$ is first- or second-order AC.

\paragraph{Scope.}
These facts identify the limits of the basic nesting dimension count,
not an impossibility theorem for every argument based on nesting.
Uniform control of the five quadrupole compressions, or stronger operator
constraints on the complement, could still lead to a better unrestricted
upper bound. No such improvement is proved here, and no catalogue upper
bound is changed by this note.
\end{document}
